% Part: turing-machines
% Chapter: undecidability
% Section: representing-tms

\documentclass[../../../include/open-logic-section]{subfiles}

\begin{document}

\olfileid{tur}{und}{rep}
\olsection{تمثيل آلات تورنغ}

\begin{explain}
كي نمثل آلات تورنغ وسلوكها بـ!!a{sentence} من منطق الرتبة الأولى، لا
بد أن نعرّف لغة مناسبة. وتتألف اللغة من جزأين: !!{predicate}s لوصف
تهيئات الآلة، وتعبيرات لترقيم خطوات التنفيذ («اللحظات») والمواضع على
الشريط.

نستحدث نوعين من ال!!{predicate}s، وكلاهما ثنائي المحل: لكل حالة~$q$،
!!a{predicate}~$\Obj Q_q$، ولكل رمز شريط~$\sigma$،
!!a{predicate}~$\Obj S_\sigma$. يتيح لنا النوع الأول وصف حالة~$M$
وموضع رأس شريطها، ويتيح لنا النوع الثاني وصف محتويات الشريط.

وللتعبير عن مواضع رأس الشريط وعدد الخطوات المنفذة نحتاج إلى طريقة
للتعبير عن الأعداد. ويتم ذلك باستعمال !!a{constant}~$\Obj 0$، ودالة
أحادية المحل~$\prime$ هي دالة الخلف. وتكتب هذه الدالة، اصطلاحًا،
\emph{بعد} وسيطها (ونحذف القوسين).

لكل عدد~$n$ حد قياسي~$\num{n}$، هو \emph{العدد القياسي} لـ~$n$، يمثله
في~$\Lang L_M$. فالحد $\num{0}$ هو $\Obj 0$، والحد $\num{1}$ هو
$\Obj 0'$، والحد $\num{2}$ هو $\Obj 0''$، وهكذا. وبصورة أدق:
\begin{align*}
\num{0} & = \Obj 0 \\
\num{n+1} &= \num{n}'
\end{align*}
يسمي الحد $\num{0}$، أي $\Obj 0$، الموضع الواقع في أقصى يسار الشريط،
وكذلك الزمن السابق لخطوة التنفيذ الأولى (التهيئة الابتدائية). ويسمي
الحد $\num{1}$، أي $\Obj 0'$، الخانة الواقعة إلى يمين الخانة الأولى،
وكذلك الزمن اللاحق لخطوة التنفيذ الأولى، وهكذا.

ونستحدث أيضًا !!a{predicate}~$<$ للتعبير عن ترتيب مواضع الشريط (حيث
يعني «إلى يسار») وترتيب خطوات التنفيذ معًا (حيث يعني «قبل»).

بعد تثبيت اللغة، ندرج «بديهيات» $!T(M,w)$، أي ال!!{sentence}s التي
تصف مجتمعةً سلوك~$M$ حين تعمل عند الدخل~$w$. وستكون هناك !!{sentence}s
تضع شروطًا على $\Obj 0$ و$\prime$ و$<$، و!!{sentence}s تصف تهيئة
الدخل، و!!{sentence}s تصف تهيئة~$M$ بعد تنفيذ تعليمة معينة.
\end{explain}

\begin{defn}
  \ollabel{defn:tm-descr}
إذا كانت $M=\tuple{Q,\Sigma,q_0,\delta}$ آلة تورنغ، فإن اللغة
$\Lang L_M$ تتألف من:
\begin{enumerate}
\item !!{predicate} ثنائي المحل $\Obj Q_q(x,y)$ لكل حالة~$q\in Q$.
  وحدسيًا تعبّر $\Obj Q_q(\num{m},\num{n})$ عن «تكون $M$، بعد $n$
  خطوة، في الحالة~$q$ وهي تمسح الخانة رقم~$m$».
\item !!{predicate} ثنائي المحل $\Obj S_\sigma(x,y)$ لكل رمز
  $\sigma\in\Sigma$. وحدسيًا تعبّر
  $\Obj S_\sigma(\num{m},\num{n})$ عن «تحتوي الخانة رقم~$m$، بعد $n$
  خطوة، الرمز~$\sigma$».
\item !!a{constant}~$\Obj 0$.
\item !!a{function} أحادية المحل~$\prime$.
\item !!a{predicate} ثنائي المحل~$<$.
\end{enumerate}
\end{defn}

ال!!{sentence}s التي تصف عمل آلة تورنغ~$M$ عند الدخل
$w=\sigma_{i_1}\dots\sigma_{i_k}$ هي الآتية:
\begin{enumerate}
\item بديهيات تصف الأعداد و~$<$:
\begin{enumerate}
\item !!^a{sentence} تقول إن كل عدد أصغر من خلفه:
\[
\lforall[x][x < x']
\]
\item !!^a{sentence} تضمن أن $<$ متعدية:
\[
\lforall[x][\lforall[y][\lforall[z][
      ((x < y \land y < z) \lif x < z)]]]
\]
\end{enumerate}
\item بديهيات تصف تهيئة الدخل:
\begin{enumerate}
\item بعد $0$ من الخطوات، أي قبل أن تبدأ الآلة، تكون $M$ في الحالة
  الابتدائية~$q_0$ وهي تمسح الخانة رقم~$1$:
\[
\Obj Q_{q_0}(\num{1},\num{0})
\]
\item تحتوي الخانات $k+1$ الأولى الرموز $\TMendtape$،
  و$\sigma_{i_1}$،~\dots، و$\sigma_{i_k}$:
\[
\Obj S_\TMendtape(\num{0}, \num{0}) \land
\Obj S_{\sigma_{i_1}}(\num{1}, \num{0}) \land
\dots \land
\Obj S_{\sigma_{i_k}}(\num{k}, \num{0})
\]
\item يكون الشريط خاليًا فيما عدا ذلك:
\[
\lforall[x][(\num{k} < x \lif \Obj S_\TMblank(x, \num{0}))]
\]
\end{enumerate}
\item بديهيات تصف الانتقال من تهيئة إلى التهيئة التالية:

فيما يأتي، لتكن $!A(x,y)$ اقتران جميع ال!!{sentence}s التي من الصورة
\[
\lforall[z][
  (((z < x \lor x < z) \land \Obj S_\sigma(z, y))
  \lif \Obj S_\sigma(z, y'))]
\]
حيث $\sigma\in\Sigma$. ونستعمل $!A(\num{m},\num{n})$ للتعبير عن
«بعد $n+1$ خطوة يكون الشريط، فيما عدا الخانة رقم~$m$، مماثلًا لما كان
عليه بعد $n$ خطوة».
\begin{enumerate}
\item \ollabel{rep-right} لكل تعليمة
$\delta(q_i,\sigma)=\tuple{q_j,\sigma',\TMright}$، ال!!{sentence}:
\begin{align*}
& \lforall[x][\lforall[y][(
   (\Obj Q_{q_i}(x, y) \land \Obj S_{\sigma}(x, y)) \lif {}]] \\
&\qquad   (\Obj Q_{q_j}(x', y') \land \Obj S_{\sigma'}(x, y') \land
!A(x, y)))
\end{align*}
تقول هذه إنه إذا كانت الآلة، بعد~$y$ خطوة، في الحالة~$q_i$ وهي تمسح
الخانة~$x$ التي تحتوي الرمز~$\sigma$، فإنها بعد $y+1$ خطوة تمسح
الخانة~$x+1$ وتكون في الحالة~$q_j$، وتحتوي الخانة~$x$ الآن الرمز
$\sigma'$، وتحتوي كل خانة غير~$x$ الرمز نفسه الذي كانت تحتويه بعد~$y$
خطوة.

\item \ollabel{rep-left} لكل تعليمة
$\delta(q_i,\sigma)=\tuple{q_j,\sigma',\TMleft}$، ال!!{sentence}:
\begin{align*}
& \lforall[x][\lforall[y][
    ((\Obj Q_{q_i}(x', y) \land \Obj S_{\sigma}(x', y)) \lif {}]]\\
& \qquad   (\Obj Q_{q_j}(x, y') \land \Obj S_{\sigma'}(x', y') \land
!A(x', y))) \land {}\\
& \lforall[y][((\Obj Q_{q_i}(\num{0}, y) \land \Obj S_{\sigma}(\num{0},
    y)) \lif {}]\\
& \qquad (\Obj Q_{q_j}(\num{0}, y') \land \Obj S_{\sigma'}(\num{0},
  y') \land !A(\num{0}, y)))
\end{align*}
توقف لحظة لتتأمل كيفية عمل ذلك: لا نبدأ الآن بعبارة «إذا كانت الآلة
تمسح الخانة~$x$،~\dots»، بل بعبارة «إذا كانت تمسح الخانة~$x+1$،
\dots». فالحركة إلى اليسار تعني أن الآلة تمسح في الخطوة التالية
الخانة~$x$. لكن الخانة التي تكتب عليها هي~$x+1$. ونفعل ذلك على هذا
النحو لأن لغتنا لا تحتوي الطرح ولا دالة السلف.

لاحظ أن الأعداد من الصورة $x+1$ هي $1$، و$2$،~\dots؛ ومن ثم لا يغطي
هذا حالة مسح الآلة الخانة~$0$ مع أمرها بالحركة يسارًا (وهو ما لا
تستطيعه بالطبع، فتبقى في موضعها). ويغطي الاقتران الثاني هذه الحالة
الخاصة: إذ يقول إنه إذا كانت الآلة، بعد $y$ خطوة، تمسح الخانة~$0$ في
الحالة~$q_i$، وكانت الخانة~$0$ تحتوي الرمز~$\sigma$، فإنها بعد $y+1$
خطوة تظل تمسح الخانة~$0$، وتصير في الحالة~$q_j$، ويصير الرمز على
الخانة~$0$ هو $\sigma'$، وتحتوي الخانات غير~$0$ الرموز نفسها التي
كانت تحتويها بعد~$y$ خطوة.
\item \ollabel{rep-stay} لكل تعليمة
$\delta(q_i,\sigma)=\tuple{q_j,\sigma',\TMstay}$، ال!!{sentence}:
\begin{align*}
& \lforall[x][\lforall[y][(
   (\Obj Q_{q_i}(x, y) \land \Obj S_{\sigma}(x, y)) \lif {}]] \\
&\qquad   (\Obj Q_{q_j}(x, y') \land \Obj S_{\sigma'}(x, y') \land
!A(x, y)))
\end{align*}
\end{enumerate}
\end{enumerate}
لتكن $!T(M,w)$ اقتران جميع ال!!{sentence}s السابقة لآلة تورنغ~$M$
والدخل~$w$.

لكي نعبر عن أن~$M$ تتوقف في نهاية المطاف، لا بد أن نجد !!{sentence}
تقول: «بعد عدد ما من الخطوات تكون دالة الانتقال غير معرّفة». ولتكن
$X$ مجموعة جميع الأزواج $\tuple{q,\sigma}$ التي تكون عندها
$\delta(q,\sigma)$ غير معرّفة. ولتكن $!E(M,w)$ عندئذ ال!!{sentence}
\[
\lexists[x][\lexists[y][(\bigvee_{\tuple{q, \sigma} \in
      X}(\Obj Q_q(x, y) \land \Obj S_\sigma(x, y)))]]
\]

وإذا استعملنا آلة تورنغ لها حالة توقف معينة~$h$ كان الأمر أيسر: عندئذ
تعبر ال!!{sentence}~$!E(M,w)$
\[
\lexists[x][\lexists[y][\Obj Q_h(x, y)]]
\]
عن أن الآلة تتوقف في نهاية المطاف.

\begin{prop}
\ollabel{prop:mlessk}
إذا كان $m<k$، فإن $!T(M,w)\Entails\num{m}<\num{k}$.
\end{prop}

\begin{proof}
تمرين.
\end{proof}

\begin{prob}
أثبت \olref[tur][und][rep]{prop:mlessk}.
(تلميح: استعمل الاستقراء على $k-m$.)
\end{prob}

\end{document}
